\section{The RMT-Cleaned Covariance Estimator}\label{sec:rmt}

\subsection{Marchenko--Pastur spectral separation}

The Marchenko--Pastur (MP) law~\cite{marchenko1967} describes the limiting spectral
distribution of $\hat\Sigma = X^\top X / T$ when $X$ has i.i.d.\ entries with
variance $\sigma^2$ and $n, T \to \infty$ with $n/T \to c \in (0,1)$. Under the
null hypothesis that $\hat\Sigma$ is a rescaled Wishart matrix (no true signal),
the bulk of the empirical spectral distribution (ESD) is supported on
$[\lambda_-, \lambda_+]$ where
\[
  \lambda_{\pm} = \sigma^2(1 \pm \sqrt{c})^2\,.
\]
In finite samples, $\sigma^2$ is estimated by the average entry-level variance
$\hat\sigma^2 = \norm{X}_F^2/(nT) = \mathrm{tr}(\hat\Sigma)/n$ (the grand mean
eigenvalue) and $c$ is replaced by $n/T$, giving the empirical edges
$\hat\lambda_\pm = \hat\sigma^2(1 \pm \sqrt{n/T})^2$.

Eigenvalues above $\hat\lambda_+$ are not explained by the null and are attributed
to genuine covariance structure. Eigenvalues below $\hat\lambda_+$ are consistent
with noise and carry no reliable directional information. We denote by $k$ the
number of signal eigenvalues: $k = |\{j : \lambda_j(\hat\Sigma) > \hat\lambda_+\}|$.
For S\&P~500 data with $n=494$ and $T=1213$, $\hat\lambda_+ \approx 1.22 \times 10^{-3}$
and $k = 23$; for the factor-calibrated synthetic model of
Section~\ref{sec:results}, the detected $k$ grows approximately as $n/21$ in this
regime.

\subsection{Construction of \texorpdfstring{$T_0^{\mathrm{RMT}}$}{T0RMT}}

Let $\hat\Sigma = U\Lambda U^\top$ be the eigendecomposition of the sample
covariance. Partition the eigenpairs by the MP edge: let
$U_k \in \mathbb{R}^{n \times k}$ hold the $k$ signal eigenvectors with
eigenvalues $\Lambda_k = \mathrm{diag}(\lambda_1, \ldots, \lambda_k)$, and let
$U_0 \in \mathbb{R}^{n \times (n-k)}$ hold the noise eigenvectors. The bulk
eigenvalues are replaced by their average,
\[
  \bar\lambda
  \;=\; \frac{1}{\,n-k\,}\sum_{j > k} \lambda_j
  \;=\; \frac{\mathrm{tr}(\hat\Sigma) - \sum_{j \leq k}\lambda_j}{n-k}\,,
\]
which preserves the trace -- total variance -- of the sample covariance.
The resolution of the identity $U_k U_k^\top + U_0 U_0^\top = I$ gives
\[
  T_0^{\mathrm{RMT}}
  = U_k \Lambda_k U_k^\top + \bar\lambda\, U_0 U_0^\top
  = \bar\lambda I + U_k \Delta_k U_k^\top\,,
\]
where
$\Delta_k = \Lambda_k - \bar\lambda I = \mathrm{diag}(\delta_1, \ldots, \delta_k)$
with $\delta_j = \lambda_j - \bar\lambda > 0$ for all $j$ (since
$\lambda_j > \hat\lambda_+ > \hat\sigma^2 > \bar\lambda$).
This is exactly the ``constant residual eigenvalue'' (CRE) -- or eigenvalue
clipping -- estimator introduced by Laloux et
al.~\cite{laloux1999}, studied in depth by Bun, Bouchaud and
Potters~\cite{bun2017}, and widely adopted in quantitative finance following
L\'opez de Prado~\cite{lopezdeprado2018}.

\begin{proposition}[Spectral properties of $T_0^{\mathrm{RMT}}$]\label{prop:rmt_spectrum}
The RMT-cleaned estimator satisfies:
\begin{enumerate}
  \item $T_0^{\mathrm{RMT}} \succ 0$ with eigenvalues
        $\{\lambda_1, \ldots, \lambda_k, \bar\lambda, \ldots, \bar\lambda\}$,
        and $\mathrm{tr}(T_0^{\mathrm{RMT}}) = \mathrm{tr}(\hat\Sigma)$.
  \item $\kappa(T_0^{\mathrm{RMT}}) = \lambda_{\max}(\hat\Sigma) / \bar\lambda$.
  \item $T_0^{\mathrm{RMT}} = \hat\Sigma - U_0(\Lambda_0 - \bar\lambda I)U_0^\top$,
        i.e., $T_0^{\mathrm{RMT}}$ differs from $\hat\Sigma$ only in the noise subspace.
  \item The matrix-vector product $T_0^{\mathrm{RMT}} v = \bar\lambda v + U_k(\Delta_k(U_k^\top v))$
        costs $O(nk)$ and requires $O(nk)$ storage.
\end{enumerate}
\end{proposition}

\begin{proof}
(1) follows from the eigendecomposition $T_0^{\mathrm{RMT}} = [U_k\;U_0]\,
\mathrm{diag}(\Lambda_k, \bar\lambda I_{n-k})\,[U_k\;U_0]^\top$,
$\delta_j > 0$ for all signal eigenpairs, and the definition of $\bar\lambda$
as the bulk average.
(2) follows from (1): the maximum eigenvalue is $\lambda_1 > 0$ and the minimum
is $\bar\lambda > 0$.
(3) is direct algebra.
(4) The product costs two dense matrix-vector products with $U_k \in \mathbb{R}^{n \times k}$.
\end{proof}

\begin{remark}[Condition number]
$T_0^{\mathrm{RMT}}$ achieves $\kappa = \lambda_{\max}(\hat\Sigma)/\bar\lambda$:
noise eigenvalues are lifted to the floor $\bar\lambda$, signal eigenvalues are
preserved. On S\&P~500 data this gives $\kappa \approx 343$, versus
$\kappa(\hat\Sigma) \approx 106{,}000$ -- an improvement by a factor of roughly
$300$. The low condition number is a consequence of replacing the
$(n-k)$-dimensional noise subspace eigenvalue spectrum (which spans
several decades) with a single value $\bar\lambda$; it also keeps the
$k \times k$ Woodbury correction matrix $W$ well-conditioned regardless of $n$.
\end{remark}

\begin{remark}[Relation to optimal RMT cleaning]\label{rem:rie}
CRE clipping is the simplest member of the family of rotationally invariant
estimators (RIE). It is not loss-optimal: the oracle RIE replaces \emph{every}
eigenvalue -- bulk and signal alike -- by a value determined by the
eigenvector overlaps of Ledoit and P\'ech\'e~\cite{ledoitpeche2011}, and these
oracle values are not constant across the bulk; nonlinear
shrinkage~\cite{ledoit2012,ledoit2020} and the RIE of Bun et
al.~\cite{bun2017} dominate clipping in Frobenius loss. In particular, CRE
retains the signal eigenvalues unshrunk although eigenvalues just above the
edge are biased upward and their eigenvectors imperfectly aligned with the
population factors (the BBP transition~\cite{baik2005,johnstone2001}). The
reason this paper nevertheless works with CRE is structural: clipping is the
member of the family whose output is \emph{low-rank-plus-identity}
(a spiked structure in the sense of Johnstone~\cite{johnstone2001}), and only
that structure admits the closed-form Woodbury solve of
Section~\ref{sec:algorithm}. Nonlinear shrinkage and general RIE produce
full-rank spectral modifications to which the Woodbury identity does not
apply; portfolios built on them require a general factorisation. Whether the
statistical gap matters for long-only minimum-variance portfolios is an
empirical question; on our data it is small
(Section~\ref{ssec:oos}).
\end{remark}

\begin{remark}[Externally specified factor models]\label{rem:external_factors}
The MP threshold provides a data-driven rule for selecting $k$, but the Woodbury
solver applies for any choice of $k$. Practitioners with an external factor model --
a commercial risk model with pre-specified style and industry factors, or a PCA
with $k$ fixed by cross-validation -- can set $k = k_{\mathrm{ext}}$ and supply
the corresponding eigenpairs directly, bypassing the MP threshold entirely. The
key structural assumption for the Woodbury solve is not the source of $k$ but
the replacement of the remaining $n - k$ eigenvalues by a single scalar $\bar\lambda$.
The rSVD preprocessing of Section~\ref{sec:preprocessing} and the complexity
bound of Proposition~\ref{prop:pipeline} hold with $k$ replaced by $k_{\mathrm{ext}}$.
When $k_{\mathrm{ext}}$ is known in advance, the preprocessing cost can be reduced
to $O(nTk_{\mathrm{ext}})$ by targeting only the required number of components
rather than applying the MP threshold post-hoc.
\end{remark}

\subsection{Choice of \texorpdfstring{$k$}{k}}\label{ssec:k_selection}

The threshold $k$ is determined by the MP upper edge $\hat\lambda_+$. In practice
$k$ is stable across the S\&P~500 sample: varying the estimation window over the
five-year period gives $k \in [21, 25]$, with the central value $k = 23$.
An audit at $k - 1$, $k$, and $k + 1$ adds negligible cost (one extra
Woodbury solve per audit point) and reveals the portfolio's sensitivity to the
threshold, which is small in the bulk of the frontier (Section~\ref{sec:ksensitivity}).

Two systematic biases in $k$ deserve mention.

\emph{Finite-sample edge fluctuations}. At finite $T$ the largest noise
eigenvalue fluctuates around the bulk edge at the Tracy--Widom scale
$O(T^{-2/3})$~\cite{johnstone2001}, so a small number of noise eigenvalues may
be classified as signal (or vice versa); refined spectrum estimators such as
El~Karoui's~\cite{elkaroui2008} reduce this effect. For the sample sizes used
here ($T \geq 250$) the impact on $k$ is at most one or two eigenvalues, which
is precisely the perturbation covered by the $k \pm 1$ audit of
Section~\ref{sec:ksensitivity}.

\emph{Sparse structure below the detection threshold}. Sector correlations below
the MP edge contribute a diffuse signal that the threshold does not detect. If
these directions are material for the portfolio, the audit at $k \pm 1$ will
show sensitivity; in practice, the minimum-variance portfolio is insensitive to
directions with small eigenvalues because those directions contribute little
variance.
